\textbf{Proof.}
Assume \(a=0\).  By the definitions of the density factors, for every \(x\in(0,1]\),
\[
 f_S(x)=x^a g_S(x)=g_S(x),
 \qquad
 f_T(x)=g_T(x),
 \qquad
 r(x)=x^{-a}\frac{g_T(x)}{g_S(x)}=\frac{g_T(x)}{g_S(x)}.
\]
The source- and target-boundary-factor assumptions give
\[
 c_{S,-}\le g_S(x)\le c_{S,+}
 \quad\text{and}\quad
 c_{T,-}\le g_T(x)\le c_{T,+},
 \qquad x\in[0,1].
\]
Because \(c_{S,-}>0\), division by \(g_S(x)\) is valid.  Thus, for every \(x\in(0,1]\),
\[
 \frac{c_{T,-}}{c_{S,+}}
 \le \frac{g_T(x)}{g_S(x)}
 =r(x)
 \le \frac{c_{T,+}}{c_{S,-}}.
\]
The restrictions \(t>1/2\) and \(s/(2s+1)+t/(2t+1)>1/2\) are compatible with these calculations but are not needed for them.

Since \(a=0<1\), the first branch in the definition of \(h_n\) yields
\[
 h_n=n^{-1/\{2(s+1)\}}.
\]
The declared condition \(s>0\) gives \(s+1>0\), so exponentiation gives
\[
 h_n^{s+1}
 =\left(n^{-1/\{2(s+1)\}}\right)^{s+1}
 =n^{-1/2}.
\]
Moreover, \(h_n\downarrow0\).  Applying the \(a<1\) branch of the Riesz-threshold lemma at \(h=h_n\) gives
\[
 \int_{h_n}^1\frac{f_T(x)^2}{f_S(x)}\,dx\asymp1,
\]
where the comparison constants depend only on the fixed source and target boundary-factor constants.  In particular,
\[
 \int_{h_n}^1\frac{f_T(x)^2}{f_S(x)}\,dx=O(1).
\]

It remains to prove the asserted positivity.  Fix a fold \(k\) for which both interior training measures are nonzero.  This condition places the construction in one of the two positive-interior-mass optimization branches of the executable-weight definition, rather than in its deterministic zero branch.  In particular, the target interior measure has positive mass, and hence
\[
 \sum_{j\in\mathcal T^T_k}\mathbf 1\{X_j^T\ge h_n\}>0.
\]
Therefore there exists \(j_0\in\mathcal T^T_k\) with \(X_{j_0}^T\ge h_n\).  Also \(|\mathcal T^T_k|>0\).  For fixed \(x\in[h_n,1]\), define
\[
 S_k(x)=\frac1{|\mathcal T^T_k|}
 \sum_{j\in\mathcal T^T_k}\mathbf 1\{X_j^T\ge h_n\}
 \exp\!\left(
 -\frac{c_{\mathrm{OT}}(x,X_j^T)
 +\widehat\theta_{h_n,J_{r,n}}^{(-k)\top}
       \widetilde{\mathbf b}_{h_n,J_{r,n}}(x)
 +\widehat\psi^{(-k)}_{h_n,J_{r,n},j}}
 {\eta_0}
 \right).
\]
The selected multipliers are finite real vectors, \(c_{\mathrm{OT}}(x,X_j^T)\) is finite, and \(\eta_0>0\), all by the executable-weight definition and the declared symbol domains.  Consequently every exponential in the sum with nonzero indicator is strictly positive.  The \(j_0\)-summand is such a term, so
\[
 S_k(x)>0.
\]
Because \(\eta_0>0\) and \(\lambda_0>0\),
\[
 \frac{\eta_0}{\eta_0+\lambda_0}>0,
\]
and the raw continuation satisfies
\[
 \widetilde r^{\mathrm{KKT},(-k)}_{h_n,J_{r,n}}(x)
 =S_k(x)^{\eta_0/(\eta_0+\lambda_0)}>0.
\]
At \(a=0\), the downstream cap is the strictly positive constant
\[
 \frac{c_{T,+}}{c_{S,-}}h_n^{-a}
 =\frac{c_{T,+}}{c_{S,-}}>0.
\]
The capped downstream weight is the minimum of this constant and the positive raw continuation.  Hence
\[
 \widehat r^{\mathrm{RICOT},(-k)}_{h_n,J_{r,n}}(x)
 =\min\!\left\{
 \widetilde r^{\mathrm{KKT},(-k)}_{h_n,J_{r,n}}(x),
 \frac{c_{T,+}}{c_{S,-}}
 \right\}>0.
\]
Since \(x\in[h_n,1]\) was arbitrary, the positivity holds throughout that interval.  The argument establishes only the stated deterministic properties; it supplies no empirical norm, realized-residual, or root-risk bound. \(\square\)